\section{Related Work}
\label{sec:related_work}

\paragraph{Preference elicitation and behavioral consequences.}
A growing literature studies whether language-model choices can be
described using preference or utility-like representations.
\citet{mazeika2025utility} model repeated choices as utility functions
and report increasingly coherent preference structures across model
scale. Other work tests whether such measured preferences correspond to
behavior outside the original elicitation setting.
\citet{tagliabue2025probing} combine verbal and behavioral preference
tests and introduce explicit costs and rewards to examine the strength
of observed behavioral tendencies. In a different setting,
\citet{slama2026preferences} find that measured preferences predict some
downstream behaviors, including donation advice and refusal patterns,
but do not consistently predict task performance.
\citet{zhou2026incentives} likewise find that coherent elicited utilities
need not translate into improved downstream performance when preferred
outcomes are attached to task success.

These findings motivate two design choices in Guardian Lens. First, we
measure behavior under an explicit resource trade-off rather than rely
only on verbal statements or cost-free choices. Second, we avoid
interpreting behavioral regularities as evidence of genuine internal
preferences. Our profiles are researcher-induced decision policies, and
our objective is to determine whether those policies can be identified
from their observable consequences.

The elicitation procedure itself can also substantially affect measured
behavior. \citet{mahajan2026mind} show that stated--revealed preference
correlations change markedly depending on whether models are forced to
choose or allowed to express neutrality and indeterminacy. This result
underscores a broader methodological point relevant to our setting:
behavioral conclusions depend on the interventions and response
constraints used to produce the observations. Guardian Lens therefore
uses fixed quantitative response schemas, matched interventions, and a
pre-specified feature representation.

\paragraph{Concept-conditioned black-box auditing.}
The closest methodological line of work concerns black-box detection of
concept-conditioned behavioral divergence. RAVEN
\citep{min2026propaganda} audits language models using two complementary
signals: within-model semantic consistency and disagreement with peer
models. Responses are clustered semantically across paraphrased prompts,
and cases with low semantic entropy but strong cross-model disagreement
receive high suspicion scores. The method is designed as a black-box
triage mechanism rather than a causal attribution tool.

VISTA \citep{liao2026vista} extends this approach to vision-language
models. It evaluates visual concept-conditioned divergence across
multiple VLMs by combining response consistency with distributional
divergence from peer models. VISTA emphasizes a central difficulty of
visual auditing: natural images contain background, composition,
lighting, and other contextual variation, so manipulating one visual
concept while preserving the surrounding scene is generally difficult.
In our deliberately controlled emblem setting, Guardian Lens addresses
this problem by constructing exact matched visual counterfactuals:
modified images are generated from the same source scene and changes are
restricted to a logged overlay region.

RAVEN and VISTA also demonstrate that black-box concept-conditioned
auditing, including visual auditing, is already an established research
direction. Guardian Lens therefore does not claim novelty from the use
of black-box queries or visual concepts alone. The distinction lies in
the identification target and experimental design. RAVEN and VISTA
primarily ask whether a model's responses are anomalous relative to a
set of peer models. Guardian Lens instead assumes known
researcher-controlled policy classes and asks whether an auditor can
recover which policy generated observed behavior.

\paragraph{Positioning of Guardian Lens.}
Guardian Lens can consequently be viewed as a controlled behavioral
system-identification benchmark rather than a semantic-divergence
detector. It differs from peer-relative auditing in four respects.
First, the model weights are fixed and the ground-truth behavioral
profile is supplied through controlled system instructions. Second,
visual effects are estimated using matched counterfactual image pairs
rather than heterogeneous images alone. Third, the output is a
quantitative allocation decision, permitting explicit measurement of
behavior under an efficiency cost. Finally, the auditor is calibrated
on separate scenes and frozen before blinded held-out profile recovery.
These choices make it possible to evaluate not only whether conditional
behavior exists, but whether distinct forms of conditional and
generalized behavior are identifiable from a pre-specified set of
black-box observations.